\expandafter\gdef\csname SNside1\endcsname{\node[anchor=north west,inner sep=0pt,text width=48mm,font=\SideFont,align=left] at ([xshift=138mm,yshift=-497.92pt]current page.north west) {\textcolor{SNink}{\textbf{错项排除 · 先决定ASL的事件空间}}\par\vspace{1.5mm}成功概率分配给已有记录，失败概率可能分配给关键字间隙或初始散列地址。先写清每项概率再计平均；不能看到有效记录有n个，就把所有失败长度都除以n。};}
\expandafter\gdef\csname SNside2\endcsname{\node[anchor=north west,inner sep=0pt,text width=48mm,font=\SideFont,align=left] at ([xshift=138mm,yshift=-475.52pt]current page.north west) {\textcolor{SNink}{\textbf{不变量 · mid比较后必须排除}}\par\vspace{1.5mm}闭区间中，相等已处理后，目标只能在mid左边或右边。若更新high=mid且区间仅两个元素，可能反复取得同一mid而不缩小；循环条件与更新式必须配套。};}
\expandafter\gdef\csname SNside3\endcsname{\node[anchor=north west,inner sep=0pt,text width=48mm,font=\SideFont,align=left] at ([xshift=138mm,yshift=-95.00pt]current page.north west) {\textcolor{SNink}{\textbf{课程连接 · 比较少不代表访存少}}\par\vspace{1.5mm}折半查找在数组上可常数时间定位中点；链表上虽能设计相同的比较决策，走到中点还要访问中间结点。算法分析要把抽象比较与实际存储访问分开。};}
\expandafter\gdef\csname SNside4\endcsname{\node[anchor=north west,inner sep=0pt,text width=48mm,font=\SideFont,align=left] at ([xshift=138mm,yshift=-95.00pt]current page.north west) {\textcolor{SNink}{\textbf{最小反例 · 孩子都合格，孙子越界}}\par\vspace{1.5mm}根10，左孩子5，5的右孩子12。每个直接孩子比较都合格，但12位于10左子树而大于10，违反全局区间；检验BST需中序或上下界递归。\par\smallskip\begin{tikzpicture}[x=1mm,y=1mm,>=Stealth,node distance=5mm]
\node[draw=SNink!50,fill=SNblue,rounded corners=1mm,text width=34mm,align=center,inner sep=1.2mm,font=\fontsize{7.5}{9}\selectfont] (m0) at (0,0) {10};
\node[draw=SNink!50,fill=SNblue,rounded corners=1mm,text width=34mm,align=center,inner sep=1.2mm,font=\fontsize{7.5}{9}\selectfont] (m1) at (0,-11) {5};
\node[draw=SNink!50,fill=SNblue,rounded corners=1mm,text width=34mm,align=center,inner sep=1.2mm,font=\fontsize{7.5}{9}\selectfont] (m2) at (0,-22) {12};
\draw[->,SNink] (m0) -- (m1);
\draw[->,SNink] (m1) -- (m2);
\end{tikzpicture}
\par{\fontsize{7.5}{9}\selectfont 10的左子树里出现12：直接父子比较仍可能全对。}};}
\expandafter\gdef\csname SNside5\endcsname{\node[anchor=north west,inner sep=0pt,text width=48mm,font=\SideFont,align=left] at ([xshift=138mm,yshift=-95.00pt]current page.north west) {\textcolor{SNink}{\textbf{错项排除 · 删除传播与插入不同}}\par\vspace{1.5mm}插入修复最小失衡子树后一般恢复原高度，更高祖先不再受本次增高影响。删除后局部修复可能继续降低高度，仍须向上；一处调整并非所有AVL更新的共同规则。};}
\expandafter\gdef\csname SNside6\endcsname{\node[anchor=north west,inner sep=0pt,text width=48mm,font=\SideFont,align=left] at ([xshift=138mm,yshift=-95.00pt]current page.north west) {\textcolor{SNink}{\textbf{不变量 · 旋转保留中序次序}}\par\vspace{1.5mm}右旋时B原右子树T的关键字介于B与A之间，转接为A左子树仍符合区间。若忘记保存T，会丢失子树；若只交换数值，结点身份和其他链接语义也可能改变。};}
\expandafter\gdef\csname SNside7\endcsname{\node[anchor=north west,inner sep=0pt,text width=48mm,font=\SideFont,align=left] at ([xshift=138mm,yshift=-95.00pt]current page.north west) {\textcolor{SNink}{\textbf{错项排除 · 红叔要先变色}}\par\vspace{1.5mm}红双亲配红叔时，旋转不是第一选择：把双亲与叔变黑、祖父变红，将冲突上移，再保证根黑。黑叔的内外侧旋转情形不能覆盖这个状态。};}
\expandafter\gdef\csname SNside8\endcsname{\node[anchor=north west,inner sep=0pt,text width=48mm,font=\SideFont,align=left] at ([xshift=138mm,yshift=-192.52pt]current page.north west) {\textcolor{SNink}{\textbf{不变量 · B树合法性有五项}}\par\vspace{1.5mm}每次操作完成后同时检查结点内有序、孩子区间、键数上下限、孩子数、叶层深度。只画出平衡外形会漏掉关键字重复或孩子接错；临时溢出框不是最终合法状态。};}
\expandafter\gdef\csname SNside9\endcsname{\node[anchor=north west,inner sep=0pt,text width=48mm,font=\SideFont,align=left] at ([xshift=138mm,yshift=-95.00pt]current page.north west) {\textcolor{SNink}{\textbf{最小反例 · 5阶不能提前随意合并}}\par\vspace{1.5mm}5阶非根至少2键、最多4键。两个合法最小兄弟各2键，加父分隔项共5键，已经溢出；删除后某侧只剩1键，才能按1+1+2=4合并。阶数约定决定合并时机。};}
\expandafter\gdef\csname SNside10\endcsname{\node[anchor=north west,inner sep=0pt,text width=48mm,font=\SideFont,align=left] at ([xshift=138mm,yshift=-95.00pt]current page.north west) {\textcolor{SNink}{\textbf{课程连接 · 索引层服务外存页访问}}\par\vspace{1.5mm}B+内部省去整条记录，通常能在同页放更多索引项，降低层数；叶链便于范围扫描。外存优势取决于页大小、缓存和记录存放，不能改写成任何输入均比内存BST快。};}
\expandafter\gdef\csname SNside11\endcsname{\node[anchor=north west,inner sep=0pt,text width=48mm,font=\SideFont,align=left] at ([xshift=138mm,yshift=-95.00pt]current page.north west) {\textcolor{SNink}{\textbf{错项排除 · 两个取模参数别混用}}\par\vspace{1.5mm}散列函数可能用小于表长的素数p，继续探查却按表长m循环。若把探查也模p，后面合法槽位永远访问不到；二次增量也应每次从原始H0计算。};}
\expandafter\gdef\csname SNside12\endcsname{\node[anchor=north west,inner sep=0pt,text width=48mm,font=\SideFont,align=left] at ([xshift=138mm,yshift=-95.00pt]current page.north west) {\textcolor{SNink}{\textbf{不变量 · 墓碑保住后续探查路径}}\par\vspace{1.5mm}已删除槽说明这里曾发生占用，后面可能仍有沿同一序列插入的元素。查找只能在从未用空槽或探查耗尽处失败；插入复用墓碑前还要排除后方已有同关键字。};}
\expandafter\gdef\csname SNside13\endcsname{\node[anchor=north west,inner sep=0pt,text width=48mm,font=\SideFont,align=left] at ([xshift=138mm,yshift=-95.00pt]current page.north west) {\textcolor{SNink}{\textbf{错项排除 · 两个取模参数别混用}}\par\vspace{1.5mm}散列函数可能用小于表长的素数p，继续探查却按表长m循环。若把探查也模p，后面合法槽位永远访问不到；二次增量也应每次从原始H0计算。};}
\expandafter\gdef\csname SNside14\endcsname{\node[anchor=north west,inner sep=0pt,text width=48mm,font=\SideFont,align=left] at ([xshift=138mm,yshift=-95.00pt]current page.north west) {\textcolor{SNink}{\textbf{不变量 · 墓碑保住后续探查路径}}\par\vspace{1.5mm}已删除槽说明这里曾发生占用，后面可能仍有沿同一序列插入的元素。查找只能在从未用空槽或探查耗尽处失败；插入复用墓碑前还要排除后方已有同关键字。};}
\expandafter\gdef\csname SNside15\endcsname{\node[anchor=north west,inner sep=0pt,text width=48mm,font=\SideFont,align=left] at ([xshift=138mm,yshift=-95.00pt]current page.north west) {\textcolor{SNink}{\textbf{课程连接 · 索引层服务外存页访问}}\par\vspace{1.5mm}B+内部省去整条记录，通常能在同页放更多索引项，降低层数；叶链便于范围扫描。外存优势取决于页大小、缓存和记录存放，不能改写成任何输入均比内存BST快。};}
\expandafter\gdef\csname SNside16\endcsname{\node[anchor=north west,inner sep=0pt,text width=48mm,font=\SideFont,align=left] at ([xshift=138mm,yshift=-95.00pt]current page.north west) {\textcolor{SNink}{\textbf{不变量 · B树合法性有五项}}\par\vspace{1.5mm}每次操作完成后同时检查结点内有序、孩子区间、键数上下限、孩子数、叶层深度。只画出平衡外形会漏掉关键字重复或孩子接错；临时溢出框不是最终合法状态。};}
\AddToShipoutPictureFG{\begin{tikzpicture}[remember picture,overlay]\ifcsname SNside\thepage\endcsname\csname SNside\thepage\endcsname\fi\end{tikzpicture}}